\chapter{Литература}
\label{ch:references}

Ниже --- работы, на которые опирается монография, в согласованном формате.
Для препринтов и внутренних материалов проекта (библиографии волтов RIN и
RIN--DTI) в конце записи в скобках указан локатор во внутренней
библиографии; это указание на местоположение источника, а не замена выходных
данных.

\begin{thebibliography}{99}
\small

\bibitem{anderson2004} J.~R. Anderson, D.~Bothell, M.~D. Byrne и др.
An Integrated Theory of the Mind. \emph{Psychological Review}, 111(4), 2004.
(Базовая активация ACT-R; см. также Anderson \& Schooler, 1991.)

\bibitem{collins1975} A.~M. Collins, E.~F. Loftus.
A Spreading-Activation Theory of Semantic Processing.
\emph{Psychological Review}, 82(6), 1975.

\bibitem{hebb1949} D.~O. Hebb. \emph{The Organization of Behavior}. Wiley, 1949.

\bibitem{hopfield1982} J.~J. Hopfield.
Neural Networks and Physical Systems with Emergent Collective Computational Abilities.
\emph{PNAS}, 79(8), 1982.

\bibitem{park2023} J.~S. Park, J.~C. O'Brien, C.~J. Cai и др.
Generative Agents: Interactive Simulacra of Human Behavior. \emph{UIST}, 2023.

\bibitem{toy2024} J.~Toy, J.~MacAdam, P.~Tabor.
Metacognition is all you need? Using Introspection in Generative Agents to Improve Goal-directed Behavior.
arXiv:2401.10910, 2024.

\bibitem{kahneman2011} D.~Kahneman. \emph{Thinking, Fast and Slow}. Macmillan, 2011.
(System~1 / System~2.)

\bibitem{nisbett1977} R.~E. Nisbett, T.~D. Wilson.
Telling More Than We Can Know: Verbal Reports on Mental Processes.
\emph{Psychological Review}, 84(3), 1977. (Конфабуляция при самообъяснении.)

\bibitem{doyle1979} J.~Doyle. A Truth Maintenance System.
\emph{Artificial Intelligence}, 12(3), 1979.

\bibitem{dekleer1986} J.~de~Kleer. An Assumption-Based TMS.
\emph{Artificial Intelligence}, 28(2), 1986.

\bibitem{quine1951} W.~V.~O. Quine. Two Dogmas of Empiricism.
\emph{The Philosophical Review}, 60(1), 1951. (Тезис недоопределённости
Дюэма–Куайна.)

\bibitem{ross1975} L.~Ross, M.~R. Lepper, M.~Hubbard.
Perseverance in Self-Perception and Social Perception.
\emph{J. Personality and Social Psychology}, 32(5), 1975. (Belief perseverance.)

\bibitem{bak1987} P.~Bak, C.~Tang, K.~Wiesenfeld.
Self-Organized Criticality. \emph{Physical Review Letters}, 59(4), 1987.

\bibitem{beggs2003} J.~M. Beggs, D.~Plenz.
Neuronal Avalanches in Neocortical Circuits. \emph{J. Neuroscience}, 23(35), 2003.

\bibitem{page1999} L.~Page, S.~Brin, R.~Motwani, T.~Winograd.
The PageRank Citation Ranking: Bringing Order to the Web. Stanford InfoLab Technical Report, 1999.

\bibitem{tong2006} H.~Tong, C.~Faloutsos, J.-Y. Pan.
Fast Random Walk with Restart and Its Applications. \emph{ICDM}, 2006.

\bibitem{spielman2019} D.~A. Spielman.
\emph{Spectral and Algebraic Graph Theory} (лекционные заметки), Yale University.

\bibitem{ortega2018} A.~Ortega, P.~Frossard, J.~Kova\v{c}evi\'c и др.
Graph Signal Processing: Overview, Challenges, and Applications.
\emph{Proceedings of the IEEE}, 106(5), 2018.

\bibitem{topping2022} J.~Topping, F.~Di~Giovanni, B.~P. Chamberlain и др.
Understanding Over-Squashing and Bottlenecks on Graphs via Curvature.
\emph{ICLR}, 2022. (Волт RIN–DTI, поз.~148.)

\bibitem{black2023} M.~Black, Z.~Wan, A.~Nayyeri, Y.~Wang.
Understanding Oversquashing in GNNs through the Lens of Effective Resistance.
\emph{ICML}, 2023. (Волт RIN–DTI, поз.~183.)

\bibitem{rusch2023} T.~K. Rusch, M.~M. Bronstein, S.~Mishra.
A Survey on Oversmoothing in Graph Neural Networks. arXiv:2303.10993, 2023.
(Волт RIN–DTI, поз.~170.)

\bibitem{bodnar2021cw} C.~Bodnar, F.~Frasca, N.~Otter и др.
Weisfeiler and Lehman Go Cellular: CW Networks. \emph{NeurIPS}, 2021.
(Волт RIN–DTI, поз.~185.)

\bibitem{bodnar2021mpsn} C.~Bodnar, F.~Frasca, Y.~Wang и др.
Weisfeiler and Lehman Go Topological: Message Passing Simplicial Networks.
\emph{ICML}, 2021. (Волт RIN–DTI, поз.~186.)

\bibitem{bodnar2022sheaf} C.~Bodnar, F.~Di~Giovanni, B.~P. Chamberlain и др.
Neural Sheaf Diffusion. \emph{NeurIPS}, 2022.

\bibitem{wei2020} R.~Wang, D.~D. Nguyen, G.-W. Wei.
Persistent Spectral Graph. \emph{Int. J. Numerical Methods in Biomedical
Engineering}, 2020. (TDA/персистентный лапласиан для биомолекул; ср. TopologyNet,
Cang \& Wei, 2017. Волт RIN–DTI, поз.~100, 103, 104, 144.)

\bibitem{occ1988} A.~Ortony, G.~L. Clore, A.~Collins.
\emph{The Cognitive Structure of Emotions}. Cambridge University Press, 1988.
(OCC-модель оценивания; ср. компонентная модель K.~Scherer.)

\bibitem{friston2010} K.~Friston.
The Free-Energy Principle: A Unified Brain Theory?
\emph{Nature Reviews Neuroscience}, 11(2), 2010. (Ошибка предсказания как сигнал значимости.)

\bibitem{speer2017} R.~Speer, J.~Chin, C.~Havasi.
ConceptNet 5.5: An Open Multilingual Graph of General Knowledge.
\emph{AAAI}, 2017. (Включая эмбеддинги Numberbatch.)

\bibitem{sap2019} M.~Sap, R.~Le~Bras, E.~Allaway и др.
ATOMIC: An Atlas of Machine Commonsense for If-Then Reasoning. \emph{AAAI}, 2019.

\bibitem{hwang2021} J.~D. Hwang, C.~Bhagavatula, R.~Le~Bras и др.
(Comet-)ATOMIC 2020: On Symbolic and Neural Commonsense Knowledge Graphs.
\emph{AAAI}, 2021.

\bibitem{west2022} P.~West, C.~Bhagavatula, J.~Hessel и др.
Symbolic Knowledge Distillation: from General Language Models to Commonsense Models.
\emph{NAACL}, 2022. (ATOMIC10x.)

\bibitem{vrandecic2014} D.~Vrande\v{c}i\'c, M.~Kr\"otzsch.
Wikidata: A Free Collaborative Knowledgebase.
\emph{Communications of the ACM}, 57(10), 2014.

\bibitem{xu2019} K.~Xu, W.~Hu, J.~Leskovec, S.~Jegelka.
How Powerful are Graph Neural Networks? \emph{ICLR}, 2019. (GIN; волт RIN–DTI,
поз.~68.)

\bibitem{corso2020} G.~Corso, L.~Cavalleri, D.~Beaini и др.
Principal Neighbourhood Aggregation for Graph Nets. \emph{NeurIPS}, 2020.
(Волт RIN–DTI, поз.~178.)

\bibitem{hu2021lora} E.~J. Hu, Y.~Shen, P.~Wallis и др.
LoRA: Low-Rank Adaptation of Large Language Models. arXiv:2106.09685, 2021.

\bibitem{rafailov2023} R.~Rafailov, A.~Sharma, E.~Mitchell и др.
Direct Preference Optimization. \emph{NeurIPS}, 2023.

\bibitem{zhou2006hg} D.~Zhou, J.~Huang, B.~Sch\"olkopf.
Learning with Hypergraphs: Clustering, Classification, and Embedding.
\emph{NeurIPS}, 2006. (Гиперграфовый лапласиан.)

\bibitem{lim2023hodge} L.-H. Lim. Hodge Laplacians on Graphs.
\emph{SIAM Review}, 62(3), 2020. (Ходж-лапласианы и высшие размерности.)

\bibitem{helamem2026} J.~Zhu, J.~Li, C.~Zhang и др.
HeLa-Mem: Hebbian Learning and Associative Memory for LLM Agents. Препринт, 2026.
(Внутренняя библиография волта RIN, поз.~53.)

\bibitem{rinvault} Библиография памяти LLM-агентов из волта RIN: APEX-MEM, TRACE,
SYNAPSE, GAAMA, MRMS, Zep/Graphiti, Persistent Sycophancy, LifeSide, LoCoMo и др.
(Внутренние материалы проекта; полные выходные данные --- в
\texttt{99 Sources/1.~Bibliography (RIN)}.)

\bibitem{rinproject} Материалы проекта RIN: модульные системные промпты
(\texttt{prompts/global\_core/*}), persona-тесты и логи Stage~0
(\texttt{logs/rin\_stage0.jsonl}), Ринциклопедии (\texttt{docs/rincyclopedia\_*}).

\end{thebibliography}
